\section{Discussion}\label{sec:discussion}

\subsection{What a transparency instrument requires that an accuracy figure does not}

A single global model predicts the next step in an administrative case about as well as the
per-cluster models do: 0.713 against 0.726 on the action alphabet. If prediction were the goal,
the clustering would be close to unnecessary.

Prediction is not the goal. What the global model cannot produce is an object. It has no
inventory of process types, no name for any of them, no representative flow a civil servant can
be shown and asked whether it matches what they do. The catalogue is attributable: each case
belongs to a named type, each type has a path, and a reader can move from a claim about
procedure to the cases that support it. That property is what the machine-readable publication
duty contemplates, and it is not measured by next-step accuracy at all.

This is why the two contributions were kept independent. Had the catalogue's value been argued
from the downstream gain, the honest 0.013 would have undermined it. The catalogue's value rests
instead on the shuffled control, which shows the grouping is real, and on regenerability, which
makes it auditable. The downstream evaluation answers a narrower question --- is per-cluster
modelling worth the effort, and where --- and it answers it partly in the negative.

\subsection{What the validity indices are actually measuring}

The indices are not broken and they are not silent. They measure partition geometry, which is a
real property, and on this corpus geometry turns out to be very nearly a restatement of how
coarse the partition is. That is the finding, and it is more useful than the null it replaced.

The evidence is a chain rather than a correlation (Section~\ref{sec:res-rq2}). Across the full
set of clusterings, silhouette rises with the share of cases a method declines to assign, and
declining to assign the hard cases raises the score on whatever remains. Restricting to the
methods that assign everything removes that channel and exposes a second: silhouette rises with
median cluster size and falls with cluster count, and cluster count is what determines how much
of the corpus clears the evaluation floor. Both channels end in \emph{how much data survives},
neither in \emph{how well the survivors are modelled}. Holding cluster size fixed, the index has
no residual association with accuracy in any size band.

The practical consequence is specific, and it is not the familiar advice to distrust the
indices. What each index encodes is a preference over granularity, and the indices do not agree
on which preference. Sweeping $k$ from 20 to 320 for the centroid clustering, silhouette peaks
at $k = 40$ on one arm and $k = 60$ on the other, in a curve that declines only weakly and
jitters by as much as the margin separating its winner from its runner-up; Calinski--Harabasz
rises monotonically across the same sweep and would select the largest $k$ offered; the elbow
criterion picks 70. Three indices, one dataset, selections an order of magnitude apart.

So we do not conclude that silhouette is the right instrument for choosing $k$, only that
choosing $k$ is the kind of decision it speaks to, and ranking partitions by expected downstream
utility is not. On this corpus a practitioner using silhouette for the latter would in effect be
ranking by coverage without knowing it. Report the indices next to a downstream measure and a
coverage figure, and the substitution stops being tempting. Our
contribution here is an instance and a mechanism rather than a theory, and it is consistent with
much better-powered work finding these indices unsuitable for choosing between clustering
approaches \citep{yerbury2025rvi} and with the older observation that a model's internal fit can
move opposite to human judgement of the same output \citep{chang2009teahaves}.

Two corrections we made to our own analysis belong here, because both were errors of exactly the
kind the section warns about. We had earlier read a reversal into this corpus --- the
best-scoring partition producing the worst paths --- and it was an artefact of computing the
consensus partitions' indices in the co-association representation rather than in the embedding
the other fourteen are scored in; it does not survive putting all sixteen on one scale. And the
apparent association between the indices and utility, once that was fixed, turned out to be the
coverage confound described above. An index computed on a population the partition selected,
scored against a metric computed on that same population, will agree with itself.

A second methodological point generalises beyond this venue. Comparing clustering algorithms on
their own evaluated populations is invalid whenever one of them can decline to assign a case.
The reversal we observed --- best on its own set, worst on a common set --- is not a small
correction; it inverts the ranking. Any comparison involving a density-based method that emits
noise should either restrict to a common set or report coverage alongside every figure.

\subsection{When per-cluster modelling earns its keep}

The evidence supports a specific and limited recommendation, and the sharpest form of it comes
from the one comparison made on a single population. For routing --- which unit handles a case
next --- per-cluster models gain 0.063 over a global model when the best family is selected, and
0.036 family-matched, consistently across cluster sizes; on large clusters they invent
essentially nothing (0.038 unseen transitions). For actions within a unit, per-cluster modelling
is not merely unprofitable but counter-productive: family-matched, on the cases both alphabets
can be scored on, it \emph{loses} 0.009 to a single global chain, in every one of the eight main
clusterings. The recommendation is therefore not ``train action models only on large clusters''
but ``do not partition for action prediction at all; pool the corpus''.

That asymmetry has a plain reading, and the sign change makes it plainer. Organisational routing
is procedure-specific: which unit receives a case encodes what kind of case it is. The actions
recorded inside a unit follow conventions of the electronic process system that are shared
across procedures, so a model pooling all cases learns them better --- it sees every procedure's
instances of the same convention, where a per-cluster model sees only its own. For an instrument
meant to show how work moves through an administration, this is a convenient result: the
representation that matters most for transparency is the one where modelling per type pays, and
the one where it does not pay is the one a transparency artefact needs least.

\subsection{Which of the sixteen is the catalogue}\label{sec:selection}

A study may report sixteen partitions; an administration cannot publish sixteen incompatible
inventories of its own procedures. Section~\ref{sec:res-rq2} makes the choice tractable in two
ways. Accuracy barely separates the clusterings, so it cannot carry the decision alone; and the
validity indices, which would ordinarily be asked to, turn out to be a restatement of
granularity rather than an independent quality signal. What is left is to select on the
quantities themselves.

We therefore restrict the candidates to the twelve clusterings that assign every case. This is
not a preference for partitional methods but the removal of a confound: a method that can
decline to assign is scored on a population it chose, and no weighting over such a mixed field
is interpretable. Table~\ref{tab:selection} scores the twelve on accuracy over the common set,
coverage, silhouette and median cluster size, min--max scaled, under four weightings.

One candidate wins all four --- the raw-arm centroid clustering whose $k$ was chosen by
silhouette, at 40 clusters of median size 116. It has the highest coverage of any full-coverage
partition (0.774), the second-highest common-set accuracy (0.6963 against a best of 0.6965), and
it wins even under the weighting that ignores the indices entirely. The two consensus partitions
come last under every weighting. We report the criterion because the alternative --- naming the
best-performing configuration from a sweep --- is the method-shopping this paper objects to
elsewhere.

That the winner is the silhouette-selected $k$ looks like a contradiction of
Section~\ref{sec:res-rq2} and is not, but it should not be read as a vindication of the index
either. Silhouette prefers coarse partitions on this data; coarse partitions leave more cases
above the evaluation floor; coverage is what our criterion rewards. The agreement is a
coincidence of direction, not evidence that the index identified a better partition --- and
Calinski--Harabasz, applied to the same sweep, prefers the opposite extreme and would have
selected a partition this criterion ranks near the bottom. We report which $k$ won and why the
criterion chose it; we do not claim the index chose it for the right reason.

%% GENERATED by `pipeline run results` - do not edit by hand.
\begin{table}[htbp]
\centering
\caption{Choosing a deliverable among the clusterings, scored on the itinerary alphabet alone.}\label{tab:selection}
\footnotesize
\begin{tabular}{lrrrrrrr}
\toprule
clustering & common t-1 & coverage & silh. & med.\ size & clusters & in clusters $\ge$150 & score \\
\midrule
exp-raw-ksil/kmeans & 0.7127 & 0.751 & 0.760 & 72 & 60 & 0.324 & \textbf{0.900} \\
exp-raw-hcoarse/hdbscan & 0.7109 & 0.759 & 0.662 & 80 & 42 & 0.631 & \textbf{0.746} \\
exp-raw-kelbow/kmeans & 0.7111 & 0.735 & 0.731 & 62 & 70 & 0.310 & \textbf{0.689} \\
exp-raw/leiden & 0.7129 & 0.727 & 0.589 & 54 & 71 & 0.331 & \textbf{0.676} \\
exp-anon-hcoarse/hdbscan & 0.7101 & 0.757 & 0.573 & 94 & 40 & 0.680 & \textbf{0.658} \\
exp-anon/kmeans & 0.7094 & 0.747 & 0.732 & 75 & 64 & 0.331 & \textbf{0.639} \\
exp-anon/hdbscan & 0.7107 & 0.745 & 0.617 & 56 & 64 & 0.363 & \textbf{0.594} \\
exp-raw/hdbscan & 0.7098 & 0.744 & 0.688 & 60 & 66 & 0.345 & \textbf{0.588} \\
exp-raw/kmeans & 0.7088 & 0.741 & 0.737 & 64 & 66 & 0.321 & \textbf{0.551} \\
exp-anon-ksil/kmeans & 0.7074 & 0.752 & 0.757 & 71 & 60 & 0.338 & \textbf{0.513} \\
exp-anon/leiden & 0.7114 & 0.711 & 0.546 & 54 & 75 & 0.267 & \textbf{0.491} \\
exp-anon-kelbow/kmeans & 0.7074 & 0.734 & 0.744 & 64 & 70 & 0.178 & \textbf{0.432} \\
exp-raw-lres/leiden & 0.7093 & 0.698 & 0.647 & 51 & 98 & 0.040 & \textbf{0.363} \\
exp-anon-lres/leiden & 0.7091 & 0.656 & 0.596 & 39 & 108 & 0.033 & \textbf{0.159} \\
\bottomrule
\end{tabular}
\par\vspace{2pt}\footnotesize\raggedright Every clustering is a candidate: no coverage filter is applied, because assigning density noise to the nearest cluster already closed the discard channel upstream. Features are min--max scaled across candidates; accuracy is scored on the common set. Scored on the 'itinerary' alphabet alone --- the alphabets are never pooled, so the deliverable is selected on one of them and that choice is a scoping decision, not a property of the data. The same candidate wins under all four weightings, including one that ignores the validity indices entirely. ``in clusters $\ge$150'' is the share of cases the partition places in clusters large enough for the low-fabrication band of Table~\ref{tab:size-trend-itinerary}.
\end{table}


Three honest qualifications. Selecting on coverage rather than accuracy is a judgement about what
a transparency artefact is for, not a result. The entry-level quality of the chosen catalogue is
uneven: on the size analysis of Section~\ref{sec:res-rq3}, 52.9\% of cases sit in clusters large
enough for their typical path to reach the quality band where fabrication is low, so nearly half
the caseload receives an entry the paper's own evidence says fabricates more freely. That share
is a property of the partition and not of the corpus --- it ranges from about a quarter to about
a half across the twelve candidates --- which is a further reason to state which partition a
catalogue came from whenever a figure like this is quoted. And the criterion is applied after the fact rather than pre-registered: we did not
commit to these weights before seeing the candidates, and a different reader might weight
interpretability, which we do not measure, above any of them. What the criterion buys is that
the choice is stated and reproducible, not that it is the only defensible one. The entries
illustrated in Table~\ref{tab:catalogue} are drawn from the partition it selects, so the
catalogue shown and the catalogue recommended are the same object.

\subsection{The record that was absent}

Returning to the premise: \citet{brauneis2018algorithmic} found that transparency requests fail
most often because the record does not exist. The catalogue described here is an attempt to
produce that record from material administrations already hold, without asking them to document
anything new.

Two properties make it a candidate for the machine-readable publication duty rather than a
research artefact. It regenerates: the pipeline rebuilds every reported figure from stored
artefacts, and a verification step fails if the document and the artefacts disagree, so the
published account cannot silently drift from the evidence. And it is bounded honestly: the
limits in Section~\ref{sec:threats} --- particularly the duration bias --- are measured and
reported rather than left for a reader to discover.

What it is not is a process model. A typical path maximises simplicity at the cost of
generalisation \citep{buijs2014quality}, and a directly-follows representation carries known
pathologies \citep{vanderaalst2019dfg}. The catalogue is a description of what recurs, offered
as a starting point for the people who run these procedures to confirm, dispute or redesign ---
which is the use the transparency principles contemplate.

\subsection{What an as-is catalogue means for accountability, and what it does not license}

The artefact has an addressee, and saying who it is belongs to the claim. In the accountability
relation of \citet{bovens2007analysing} the actor is the administration that runs the
procedures, and the forum is plural: the internal control and ombudsman units that receive
users' complaints, the external body that audits administrative performance, and the units
charged with redesigning procedures. The service-charter duty in art.~7, \S~2, III
\citep{brasil2017cdu} adds the user of the service and fixes what the account must contain,
namely the main steps by which a service is processed. A catalogue derived from the log supplies
that content in a form that regenerates and that attributes: a statement about a procedure can
be traced to the cases behind it, which is the difference between releasing an event log and
giving an account of one \citep{ananny2018seeing, janssen2012benefits}.

The scope limits in Section~\ref{sec:legal} bind that use. One of them deserves restating here
because it is the limit a forum will most want to cross: this is not a duration assessment.
The corpus cannot support one (Section~\ref{sec:threats}), and a principle satisfied in degrees
would not yield a breach finding even from a corpus that could.

Which of the four the artefact actually serves today should be said rather than left to the
reader. Redesign units are served: an inventory of what recurs, with a representative flow per
type, is directly usable input to procedure redesign, and it is the use every scope limit
survives. An external auditor is served for inventory only, since no duration or efficiency
statistic from this corpus is admissible. Internal control and ombudsman units are \emph{not}
served in the way the framing suggests, because they act on individual complaints and this paper
forbids per-case use of an entry. And publication to the service user is not yet defensible:
a path that matches under half of a real case's visits, drawn from a corpus that excludes the
slow cases, with an unvalidated generated name, would understate what a user must go through and
would understate it most where the understatement misleads. One forum of four, conditionally, is
the honest count.

\subsection{Legibility, and the record as a thing to be managed}

Making an administration's operation legible from its own records is not a neutral act, and the
standing objection is \citet{scott1998seeing}: schemes that render practice legible to a centre
simplify it, and the simplification then acts back on the practice it claimed only to describe.
Our own measurements are what give that objection force here rather than deflecting it. The
generated paths cover fewer than half of a real case's visits, so the artefact simplifies, and
in a measured direction. Constrained decoding enforces a movement rule that observed practice
breaks about 8\% of the time, so the published account is normalised relative to the conduct it
describes. And the artefact is offered to the units that redesign procedures, which is precisely
the mechanism by which a description becomes a template.

Three consequences follow for anyone adopting it, and they are design constraints rather than
caveats. A catalogue entry must be labelled descriptive and partial, with its case count and its
support measures travelling beside it, so that a weak entry cannot be read as a strong one.
It must not be used as a conformance standard: a case that departs from its type's typical path
has not thereby deviated from anything, because the path is a summary of what recurred and not a
rule. And the classification itself becomes an object of management once published --- what
\citet{espeland2007rankings} call reactivity --- so the types should be expected to shape the
practice they describe, and be revised against the log rather than defended. The disclaimer in
Section~\ref{sec:legal} that these types are descriptive and revisable is the right instinct;
this is the literature explaining why that disclaimer erodes unless the revision is built in.

\subsection{What adoption would require}\label{sec:adoption}

The eleven-hour figure quoted in Section~\ref{sec:corpus} is the cost of the \emph{study}, not of
the artefact, and conflating them overstates what adoption costs by an order of magnitude. That
figure covers sixteen clusterings across two text arms, both trace alphabets, five model families
per cluster and forty evaluation cells --- 8{,}985 trained models whose purpose was to establish
which parts of the method are load-bearing. An administration producing a catalogue does not
repeat that. It needs ingestion into an event knowledge graph, the feature blocks, \emph{one}
clustering, the naming step, \emph{one} path learner, and the constrained decoder. On the
evidence here it should choose the routing alphabet; the model family is a cost decision the
paper reports on both sides rather than settling. The strongest family on routing is the GRU,
gaining 0.053 over the global model against the third-order Markov chain's 0.035 --- half as
much again, and not a rounding difference. Markov is the recommendation only where the GPU
infrastructure the neural families require is not already in place, at the cost of roughly a
third of the gain. On actions per-cluster modelling should not be done at all. That configuration is a small fraction
of the study's compute and needs no training hardware beyond what the document embedding
requires.

What does not shrink is everything around the models, and it is the harder half. The method
presumes an electronic process system whose event log is complete enough to reconstruct
per-case routing; a graph database and the engineering to maintain it; a document store from
which opening documents can be retrieved and parsed; and locally executed embedding and
language models, which in this study was a legal requirement rather than a preference and
therefore sets a hardware floor an administration cannot avoid by using a hosted service. That
floor was met here without procurement: the embedding and language models ran on GPU
infrastructure the state administration already owned and operates, reached over an
OpenAI-compatible interface that never leaves its network. The requirement is therefore
demanding but not exotic --- an administration that has already built shared computing capacity
can satisfy it, and one that has not must treat that, and not the modelling, as the first cost.
The method also presumes an instrument authorising the processing, which took longer to obtain
than any computation reported here.

Two consequences follow that bear on who can use this. Scale matters, because the method is
useless below the point where types have enough cases to model: at the size threshold observed
here, an administration handling a few hundred cases a year would populate no entry at all, and
the catalogue is an instrument for bodies with caseloads in the thousands. And the artefact must
be refreshed, since a published account of processing steps that silently ages is worse than
none; the regenerability described in Section~\ref{sec:setup} is what makes a refresh cheap, but
someone must own the cadence, and entries that appear or vanish between refreshes are a
governance question we do not answer.

\subsection{Governance of the unit-level record}\label{sec:governance}

The artefact is a behavioural record of two populations, and only one of them receives any
protective treatment in the method. The parties to the cases are covered by the de-identification
of Section~\ref{sec:anon}. The public servants who processed those cases are not: the routing
alphabet is, by construction, a record of which organisational unit handled what and in what
order, and it is the representation this paper argues is the more informative of the two. The
paper's best result is therefore also its most sensitive artefact, and a disclaimer that the
catalogue is not a basis for action against identified staff binds the present authors and
nobody who later inherits the pipeline.

Four requirements follow for an administration adopting this, and we state them as requirements
because the study itself did not have to meet them and a deployment would. First, unit-level
statistics approach individual-level statistics as unit headcount falls, and leaf units in a
secretariat are frequently staffed by one to three people; a published figure attributed to such
a unit is a productivity record for named individuals. A headcount or case-volume threshold
below which no unit-level figure is published is the direct analogue of the 30-trace evaluation
floor this study already applies for a statistical reason. Second, the artefact should be
separated into a published catalogue, an internal catalogue carrying full unit attribution under
a stated purpose limitation, and the underlying graph, with distinct access regimes and a named
decision-maker for each. Third, the foreseeable misuses --- unit ranking, individual performance
assessment, disciplinary use --- should be named in the instrument authorising the work rather
than left to be discovered, since nothing technical in the method prevents them. Fourth, a
per-unit routing catalogue is a matter on which staff representative bodies have a legitimate
interest; this study did not consult them, and an administration deploying it should.

We flag one further gap rather than fill it. The legal basis on which a public body processes
these records for analysis of this kind is a question of data-protection law that
Section~\ref{sec:legal} does not resolve, because this study proceeded under a controller's
authorisation and had no occasion to identify the statutory hypothesis an adopting body would
need to rely on. That question would be the first one asked of any administration attempting to
reproduce this work, and it deserves a treatment we are not in a position to give it.
